\documentclass[11pt]{article}
\usepackage{amsmath,amsthm,amssymb}
\newtheorem{theorem}{Theorem}[section]
\newtheorem{definition}[theorem]{Definition}
\newtheorem{lemma}[theorem]{Lemma}
\newtheorem{corollary}[theorem]{Corollary}
\newtheorem{proposition}[theorem]{Proposition}
\title{Mathematical Foundations of Process Scheduling, Memory Management,
       and Resource Lifecycle}
\author{Extracted from Linux 0.11 source, Oct 2026}
\date{}
\begin{document}
\maketitle

% ============================================================
\section{Process State Automaton}
% Source: sched.c states, exit.c transitions, fork.c creation
% ============================================================

\begin{definition}[State Space]
Let $S = \{R, I, U, Z, T\}$ be five states:
Running, Interruptible, Uninterruptible, Zombie, Terminated.
A process is a tuple $p = (id, s, c, \pi, \sigma, \beta, \alpha, f)$ where
$s \in S$, counter $c \in \mathbb{Z}$, priority $\pi \in \mathbb{N}^+$,
signal bitmap $\sigma \in \{0,1\}^{32}$, blocked mask $\beta \in \{0,1\}^{32}$,
alarm deadline $\alpha \in \mathbb{N} \cup \{0\}$, and father $f \in \mathbb{N}$.
\end{definition}

\begin{definition}[Blockable Mask]
Signals $k$ and $t$ are unblockable. Define
$B = \{0,1\}^{32} \setminus \{e_k, e_t\}$
where $e_i$ is the $i$-th standard basis vector.
A process's effective pending set is $\sigma \wedge \lnot(\beta \wedge B)$.
\end{definition}

\begin{theorem}[State Transition Rules]
The automaton has exactly these transitions:
\begin{enumerate}
\item $R \to I$: voluntary yield (pause).
\item $R \to U$: sleep on uninterruptible wait queue.
\item $I \to R$: wake-up when effective pending $\neq \mathbf{0}$.
\item $U \to R$: explicit wake-up by the resource holder.
\item $R \to Z$: exit; release address space, close files, notify father.
\item $Z \to \emptyset$: father collects; page freed.
\item $\forall s$: alarm fires when global tick $j \geq \alpha$; sets signal bit $a$.
\end{enumerate}
No other transitions exist. $Z$ is absorbing until collected.
\end{theorem}


% ============================================================
\section{Scheduling: Priority Decay with Starvation Freedom}
% Source: sched.c schedule(), lines 104--142
% ============================================================

\begin{definition}[Counter Recalculation]
When all runnable counters are zero, every process $p_i$ (including sleeping)
receives:
\[
  c_i' = \lfloor c_i / 2 \rfloor + \pi_i
\]
\end{definition}

\begin{theorem}[Bounded Counter]
After any number of recalculations without consumption,
$c_i < 2\pi_i$. Proof: let $c_i^{(n)}$ be the counter after $n$
recalculations without decrement.
$c_i^{(0)} = 0$.
$c_i^{(n+1)} = \lfloor c_i^{(n)}/2 \rfloor + \pi_i$.
This is a contraction mapping on $[0, 2\pi_i)$ converging to
$c^* = 2\pi_i - 1$ from below. Therefore $c_i < 2\pi_i$ always.
\end{theorem}

\begin{theorem}[Starvation Freedom]
Every runnable process is scheduled within $N \cdot 2\pi_{\max}$ ticks,
where $N = |$runnable$|$. Proof: each tick decrements the current process's
counter. After at most $2\pi_{\max}$ ticks, the current process reaches
$c = 0$. If all $N$ processes have $c = 0$, recalculation fires and the
highest $\pi$ runs. Since $\pi > 0$ for all processes, every process gets
$c > 0$ after recalculation, so it will eventually be the maximum.
\end{theorem}

\begin{corollary}[Interactive Bias]
A sleeping process accumulates credit:
after $k$ recalculations while sleeping,
$c = \pi(1 + 1/2 + \ldots + 1/2^{k-1}) < 2\pi$.
Upon waking it has higher counter than a CPU-bound process
that was consuming its counter. This gives I/O-bound processes
implicit priority boost without changing $\pi$.
\end{corollary}


% ============================================================
\section{Timer Queue: Differential Encoding}
% Source: sched.c add_timer(), do_timer(), lines 264--336
% ============================================================

\begin{definition}[Differential Timer Chain]
A timer chain is a singly-linked list $t_1 \to t_2 \to \ldots \to t_n$
where each node stores a \emph{differential delay} $d_i \geq 0$.
The absolute deadline of $t_k$ is $\sum_{i=1}^{k} d_i$.
\end{definition}

\begin{theorem}[O(1) Tick Processing]
Each global tick decrements only $d_1$. When $d_1 \leq 0$, fire $t_1$
and advance the head. No other nodes are touched.
Cost per tick: $O(1)$ amortized (each fired timer is processed once).
\end{theorem}

\begin{proposition}[Insertion by Bubble Sort]
A new timer with absolute deadline $D$ is inserted at the head
with $d = D$, then bubbled toward the tail by swapping payloads
(not pointers) with its successor whenever $d_{\text{next}} < d_{\text{current}}$.
Each swap subtracts $d_{\text{next}}$ from $d_{\text{current}}$ and adds it to
$d_{\text{next}}$, preserving all absolute deadlines.
Worst-case $O(n)$ per insert, but $n \leq 64$ (static pool).
\end{proposition}


% ============================================================
\section{Sleep Queue: Implicit Stack Discipline}
% Source: sched.c sleep_on(), interruptible_sleep_on(), wake_up()
% ============================================================

\begin{definition}[Wait Channel]
A wait channel is a pointer-to-pointer $\texttt{*p}$.
$*p$ points to the most recently sleeping process on that channel.
Each sleeper saves $\text{tmp} = *p$, then sets $*p = \text{self}$.
This forms an implicit LIFO stack in the local variables of sleeping processes.
\end{definition}

\begin{theorem}[Cascading Wake-up]
When $\text{wake\_up}(p)$ sets $(*p).\text{state} = R$ and $*p = \text{nil}$,
it wakes only the most recent sleeper. That sleeper, upon resuming from
$\text{schedule}()$, checks its saved $\text{tmp}$ and wakes that process too.
The chain unwinds in LIFO order, waking all sleepers.
\end{theorem}

\begin{lemma}[Interruptible Variant Correctness]
In the interruptible variant, a spuriously woken process (one that was
not at the head when signal arrived) must re-sleep after waking the
current head. The goto-repeat loop ensures that only the head
actually proceeds; all others re-enter sleep. This prevents
lost-wake-up bugs.
\end{lemma}


% ============================================================
\section{Process Creation: Fork as Structure Morphism}
% Source: fork.c copy_process(), copy_mem()
% ============================================================

\begin{definition}[Fork Morphism]
Fork is a function $\phi: P \to P$ where $P$ is the process space.
Given parent $p$, child $\phi(p)$ satisfies:
\begin{align}
\phi(p).\text{id}   &= \text{next\_pid}() \\
\phi(p).\text{father} &= p.\text{id} \\
\phi(p).c            &= \phi(p).\pi  \quad\text{(fresh quantum)} \\
\phi(p).\sigma       &= \mathbf{0}   \quad\text{(no pending signals)} \\
\phi(p).\alpha       &= 0            \quad\text{(no alarm)} \\
\phi(p).\text{leader} &= 0           \quad\text{(leadership doesn't inherit)} \\
\phi(p).\text{times}  &= \mathbf{0}  \quad\text{(accounting reset)}
\end{align}
All other fields (file descriptors, root, cwd, executable) are
\emph{shared by reference}: their reference counts are incremented,
not copied.
\end{definition}

\begin{theorem}[Address Space Partitioning]
Each process $n$ receives linear base $n \cdot 64\text{MB}$.
The page tables are copied with all pages marked read-only
in both parent and child (copy-on-write setup).
The physical pages are shared; $\text{mem\_map}[f]$ incremented for each.
\end{theorem}

\begin{definition}[PID Allocation]
$\text{next\_pid}()$ increments a global counter $\ell$ (wrapping at overflow
to 1), rejecting any value already in the task table.
This is a linear probe over a sparse set: expected $O(1)$ when
$|\text{tasks}| \ll \text{pid\_max}$.
\end{definition}


% ============================================================
\section{Process Termination: Exit and the Orphan Invariant}
% Source: exit.c do_exit(), sys_waitpid(), tell_father()
% ============================================================

\begin{definition}[Orphan Reparenting]
On exit of process $p$, for every process $q$ with $q.f = p.\text{id}$,
set $q.f = 1$ (init). If $q.s = Z$, signal init with SIGCHLD.
This maintains the invariant: every non-init process has a living father.
\end{definition}

\begin{theorem}[Resource Release Ordering]
Exit releases in this exact order:
(1) address space pages, (2) open files, (3) cwd inode,
(4) root inode, (5) executable inode, (6) controlling terminal.
Then state transitions to $Z$ and father is signaled.
The ordering ensures no resource is used after release:
files may reference inodes, so files close before inodes release.
\end{theorem}

\begin{theorem}[Wait Filter Semantics]
$\text{waitpid}(w, \ldots)$ scans children $q$ of current:
\begin{itemize}
\item $w > 0$: match $q.\text{id} = w$ (specific child)
\item $w = 0$: match $q.\text{pgrp} = \text{current}.\text{pgrp}$ (same group)
\item $w = -1$: match any child
\item $w < -1$: match $q.\text{pgrp} = -w$ (specific group)
\end{itemize}
If a zombie is found, accumulate its times into parent, release it,
return. If living children exist but none zombie, sleep on SIGCHLD.
If no children match at all, return ECHILD.
\end{theorem}


% ============================================================
\section{Memory: Two-Level Page Table and Copy-on-Write}
% Source: mm/memory.c
% ============================================================

\begin{definition}[Page Table Structure]
Virtual address $v$ decomposes as $(d, t, o)$ where
$d = v[31:22]$ (directory index, 10 bits),
$t = v[21:12]$ (table index, 10 bits),
$o = v[11:0]$ (offset, 12 bits).
Physical address: $\text{pg\_dir}[d][t].\text{frame} \cdot 4096 + o$.
\end{definition}

\begin{definition}[Reference-Counted Physical Pages]
$\text{mem\_map}: \{0, \ldots, P-1\} \to \mathbb{N}$
where $P = \lfloor 15\text{MB} / 4096 \rfloor$.
$\text{mem\_map}[f] = 0$ means free, $\geq 1$ means in use.
$\text{get\_free\_page}()$ scans from high address down for the first $0$ entry.
\end{definition}

\begin{theorem}[Copy-on-Write Correctness]
After fork, both parent and child page table entries point to
the same physical frame with write bit cleared. On write fault:
\begin{enumerate}
\item If $\text{mem\_map}[f] = 1$ (sole owner): set write bit, return.
      Cost: $O(1)$, no copy.
\item If $\text{mem\_map}[f] > 1$ (shared): allocate new frame $f'$,
      copy 4096 bytes, decrement $\text{mem\_map}[f]$, point PTE to $f'$.
      Cost: $O(\text{page\_size})$.
\end{enumerate}
Invariant: after handling, the faulting process has a private writable page.
\end{theorem}

\begin{theorem}[Page Sharing via Executable Identity]
Two processes sharing the same executable inode (detected by
$\text{inode.i\_count} \geq 2$) can share clean code pages.
For a page fault at virtual address $a$:
if any other process with the same executable has a clean, present
page at offset $a$, mark it read-only in both, increment
$\text{mem\_map}$, and return without I/O.
This amortizes disk reads across $n$ instances of the same binary.
\end{theorem}


% ============================================================
\section{Bucket Allocator: Size-Class Freelists}
% Source: lib/malloc.c
% ============================================================

\begin{definition}[Size Classes]
$\mathcal{C} = \{16, 32, 64, 128, 256, 512, 1024, 2048, 4096\}$.
For request size $r$, select $c = \min\{s \in \mathcal{C} : s \geq r\}$.
Internal fragmentation $\leq c - r < c$.
Worst-case waste ratio: $< 50\%$ (request $c/2 + 1$ gets bucket $c$).
\end{definition}

\begin{definition}[Bucket Page]
A bucket page of class $c$ contains $\lfloor 4096/c \rfloor$ objects
chained as a freelist. A bucket descriptor $(page, next, freeptr, refcnt, size)$
tracks the page. When $refcnt = 0$, the entire page returns to the system.
\end{definition}

\begin{theorem}[Allocation Complexity]
$\text{malloc}(r)$: scan size classes $O(|\mathcal{C}|) = O(1)$ (constant 9),
then walk bucket chain for a non-empty freeptr: $O(b)$ where
$b$ is the number of full pages of that class. Pop from freelist: $O(1)$.
Expected: $O(1)$ when pages are not exhausted.

$\text{free}(obj)$: find bucket by page address $O(|\mathcal{C}| \cdot b)$
worst case but constant in practice. Push to freelist: $O(1)$.
If $refcnt$ drops to $0$: unlink bucket descriptor, return page.
\end{theorem}


% ============================================================
\section{Signal Delivery: Stack Frame Surgery}
% Source: kernel/signal.c do_signal()
% ============================================================

\begin{definition}[Signal Disposition]
For signal $n$, disposition $h_n \in \{\text{SIG\_DFL}, \text{SIG\_IGN}, f\}$
where $f$ is a user-space handler address. Flags:
ONESHOT (reset to DFL after delivery), NOMASK (don't block self during handler).
Blocked mask $\beta$: if bit $n$ set and $n \neq k$ (SIGKILL), signal is deferred.
\end{definition}

\begin{theorem}[User Stack Frame Construction]
Signal delivery constructs a return frame on the user stack:
\[
  \text{esp} \leftarrow \text{esp} - L \cdot 4
\]
where $L = 7$ (NOMASK) or $L = 8$ (with blocked mask save).
The frame contains, bottom to top: restorer address, signal number,
(optionally) old blocked mask, saved eax, ecx, edx, eflags, old eip.
The kernel return EIP is overwritten with $f$, so on return-to-user
the handler executes. The restorer pops the frame and calls sigreturn.
This is a homeomorphism: the user stack state is recoverable.
\end{theorem}


% ============================================================
\section{Initialization as Resource Binding}
% Source: init/main.c, sched_init(), mem_init()
% ============================================================

\begin{theorem}[Boot Sequence as Monotone Binding]
Initialization is a sequence of binding functions,
each mapping an unknown resource extent to a known data structure:
\begin{enumerate}
\item $\text{mem\_init}(start, end)$: binds physical pages to $\text{mem\_map}$.
\item $\text{sched\_init}()$: binds task[0] to GDT entries, timer to IRQ 0.
\item $\text{hd\_init}()$: binds disk geometry to request queue.
\item $\text{tty\_init}()$: binds serial/console to tty structures.
\end{enumerate}
Each binding is monotone (never undone) and must precede any use
of the bound resource. The dependency order is a DAG:
sched depends on mem (task[0] needs a page), fs depends on hd
(mount needs disk), etc.
\end{theorem}

\begin{corollary}[Reference Repository Initialization]
By analogy, a system that consults reference code repositories
binds them at boot, not on demand. The binding maps
(owner, name) to a local clone directory, exactly as
$\text{mem\_init}$ maps physical address ranges to $\text{mem\_map}$ entries.
Mathematical theorems are extracted from the bound repositories
by a conversational process (not an automaton), just as
device probing discovers hardware capabilities that no
static table can enumerate.
\end{corollary}

\end{document}
